\begin{lemma}[Stable common-missing colors are closed on the PE core]
Let $T$ be a positive-rung extended Tashkinov tree constructed from a
$k$-triple $G,e_0,\varphi$, with arbitrary certified input coloring,
terminal ladder tree $T_n$, terminal generating coloring $\varphi_n$,
and terminal stored-color set $D_n$.  Assume that $T$ has maximum
property and that its terminal extension is a PE.  Let $R_n$ and
$T_n\vee R_n$ be the closure and exact ordered union in a hierarchy
base under $\varphi_n$ (respectively `base.closure' and `base.core'
in Lean).

Assume, uniformly over every certified input coloring and generating
coloring, that every positive-rung ETT $T'$ of rung number $s<n$ with
maximum property satisfies all six lower-rung assertions:
\begin{enumerate}
\item $V(T')$ is elementary under its generating coloring;
\item under every proper coloring leaving exactly $e_0$ uncolored and
stable on its final ladder closure, that closure is interchangeable for
two colors whenever at least one of them is missing there;
\item stability descends at each PE support to the prefix before the
support and to its first-appearance ladder stage, and stored pairs at
distinct PE slots with distinct supports and the same first-appearance
index are disjoint;
\item at a terminal PE, the stored two-color component through the
support meets the terminal ladder tree only at that support;
\item every nonminimum endpoint of a defective boundary color precedes
each later PE support or SE extension vertex allowed by its
first-appearance index; and
\item every proper terminal stable coloring leaving exactly $e_0$
uncolored is a mod-coloring of the ETT.
\end{enumerate}

Let $\sigma_n$ be proper, leave exactly $e_0$ uncolored, and be
$T_n\oplus R_n,D_n,\varphi_n$-stable.

For every color $c\in[k]$, if
$c\in\overline\sigma_n(T_n)\cap\overline\sigma_n(R_n)$, then
$\{c\}$ is closed on $T_n\vee R_n$ under $\sigma_n$.
All boundaries, trees, and two-color components use the original edge
type, so distinct parallel edges retain their identities.
\end{lemma}
\begin{proof}
Write $X=V(T_n)$, $Y=V(R_n)$, and
$U=V(T_n\vee R_n)$, with vertex sets understood as in
\noderef{TreeSequenceVertexSet}.  The generating-color field of
\noderef{ExtendedTashkinovTree} identifies $\varphi_n$ with the last
coloring of the certified series.  Eliminate the last certificate in
\noderef{TashkinovSeries} using the terminal-PE hypothesis.  The PE
constructor of \noderef{TashkinovSeriesStep} supplies the exact support
$v_n$, exact connecting-edge identity $f_n$, and colors
$\gamma_n,\delta_n$.  It also supplies properness and exact
one-uncoloredness of $\varphi_n$, and the equations recording the
terminal $\gamma_n,\delta_n$-switch.  No endpoint pair is substituted
for $f_n$.

The parallel clause of \noderef{HierarchyBase} says that $R_n$ is a
$\varphi_n$-closure of the ordered prefix $T_n(v_n)$, and supplies
an \noderef{OrderedTreeUnion} certificate with
\[
                         U=X\cup Y.                       \tag{1}
\]
It also says that the missing-color set of $T_n$ under the base
coloring is its missing-color set under $\varphi_n$.
The two conjuncts of \noderef{CombinedStableColoring}, with stability
expanded as in \noderef{StableColoring}, say that
$\sigma_n$ is $(T_n,D_n,\varphi_n)$-stable and simultaneously
$(R_n,\varnothing,\varphi_n)$-stable.  Applying
\noderef{CombinedStableClosurePreserved} gives, for every displayed edge
identity $g$ of either tree,
\[
             \sigma_n(g)=\varphi_n(g),                    \tag{2}
\]
and proves that the same ordered tree $R_n$, with the same edge
identities, is a $\sigma_n$-closure of $T_n(v_n)$.

We record three preliminary facts, with their hypotheses checked at the
coloring at which they will be used.

First, $U$ is elementary under $\sigma_n$.  Indeed, retain the same
core, closure, and support position from the hierarchy base.  Equality
of missing-color sets on $T_n$ is the vertex clause of the first stable
relation.  The preceding paragraph supplies the required
$\sigma_n$-closure, and the old ordered-union certificate is unchanged.
These data form a hierarchy base under $\sigma_n$.  The first stable
relation, properness, exact one-uncoloredness, positivity, maximum
property, the terminal-PE equation, and the displayed six lower-rung
assertions are precisely the hypotheses of
\noderef{ParallelClosureUnionElementary}.  Hence
\[
                  U\text{ is elementary under }\sigma_n.   \tag{3}
\]

Second, every color missing on both $T_n$ and $R_n$ is closed on
their union under $\sigma_n$.  For the pre-step elementarity premise of
\noderef{ParallelCommonMissingClosed}, distinguish the two permitted
positive-rung cases.  If $n>1$, apply
\noderef{CertifiedETTPredecessor}.  It returns an actual certified ETT
of rung $n-1$, with maximum property, terminal tree exactly $T_n$,
generating coloring exactly $\varphi_{n-1}$, and entrywise equality of
all retained arrays and certificates.  Its rung is positive and smaller
than $n$, so lower assertion (i), instantiated with its displayed input
and generating coloring, proves that $V(T_n)$ is elementary under
$\varphi_{n-1}$.  If $n=1$, slot zero is the initial tuple and the same
claim is the first conclusion of \noderef{TechnicalTheoremBase}.
Thus every hypothesis of \noderef{ParallelCommonMissingClosed} is now
verified, and that lemma gives
\[
 T_n\vee R_n\text{ is }
 \bigl(\overline\varphi_n(T_n)\cap
       \overline\varphi_n(R_n)\bigr)\text{-closed under }\varphi_n.
                                                               \tag{4}
\]
The vertex clauses of the two stable relations identify both constituent
missing-color sets under $\varphi_n$ and $\sigma_n$.  The PE constructor
of \noderef{TashkinovSeriesStep} says that $\varphi_n$ is proper and
leaves exactly $e_0$ uncolored; the same two facts for $\sigma_n$ are
hypotheses.  We may therefore apply the symmetry clause of \noderef{sc2}
separately to the two constituent stable relations.  It gives the reverse
relations saying that $\varphi_n$ is $(T_n,D_n,\sigma_n)$-stable and is
also $(R_n,\varnothing,\sigma_n)$-stable.

Now let $c$ be common missing under $\sigma_n$, and suppose that an exact
edge identity $g$ lies in the $c$-colored boundary of $U$ under
$\sigma_n$.  By \noderef{BoundaryEdges}, name its endpoints so that the
inside endpoint $p$ lies in $U$ and the other lies outside $U$.  Equation
(1) gives $p\in X$ or $p\in Y$, and \noderef{TreeIncidentEdge} makes $g$
incident with the corresponding constituent tree.  If $p\in X$, then
$c\in\overline\sigma_n(T_n)$ and the edge clause of the reverse
$T_n$-stable relation gives $\varphi_n(g)=\sigma_n(g)=c$.  If $p\in Y$,
the reverse $R_n$-stable relation gives the same equality.  Boundary
membership depends only on the exact edge endpoints, while the color
filter is \noderef{BoundaryColorEdges}; hence $g$ is a $c$-colored
boundary identity of $U$ under $\varphi_n$.  The vertex clauses also put
$c$ in both constituent missing-color sets under $\varphi_n$, so this
contradicts (4).  Consequently
\[
 c\in\overline\sigma_n(T_n)\cap\overline\sigma_n(R_n)
 \quad\Longrightarrow\quad
 \partial_{\sigma_n,c}(U)=\varnothing.                       \tag{5}
\]
This is a statement about exact members of the boundary set in
\noderef{TreeColorSetClosed}; common missing membership has its vertex
witness by \noderef{TreeMissingColors}.
\end{proof}
